This study derives direction-dependent characteristic lengths from ellipsoidal nonlocal averaging and incorporates them into a unified finite element framework for anisotropic strain-gradient elasticity. The anisotropic length tensor $L_{mn}=(AA^T)_{mn}$ is obtained as second moments of the averaging ellipsoid, giving a geometric interpretation analogous to fabric tensor. A positive-definite second-order strain-gradient energy $W=\frac12 C_{ijkl}\varepsilon_{ij}\varepsilon_{kl}+\frac1{20}L_{mn}C_{ijkl}\varepsilon_{ij,m}\varepsilon_{kl,n}$ is formulated and implemented with $C^1$-continuous cubic Hermite (1D), Bogner-Fox-Schmit bicubic Hermite (2D) and tricubic Hermite (3D, 24 DOF/node) elements. Static deformation, free vibration, anisotropic length effects, orientation-angle effects and AR-$\theta$ design-space ($AR=l_x/l_y\in[1,20]$, $\theta\in[0^\circ,90^\circ]$) are explored. With corrected consistent mass $M_{scaled}=S_{scale}M_{ff}S_{scale}$, the orientation sweep for $l_x=0.20$ m, $l_y=l_z=0.02$ m shows monotonic variation: $u(\theta)/u(0^\circ)=1.00\to1.039$, $f_1=186.26\to173.06$ Hz, $f_{iso}=178.72$ Hz for $l_{ref}=0.05$ m ($f_{iso}=177.0$ Hz for $l_{iso}=0.0431$ m), $f_1/f_{iso}=1.054\to0.979$, $\kappa=2.7\times10^5\to10.3\times10^5$, $u_{ref}=4.7876$ $\mu$m, satisfying $f_1>171.61$ Hz classical bound. Axial length dominates transverse ($5.8\%$ vs $0.5\%$ reduction). DSI=$(u_{ref}-u)/u_{ref}$ with volume-equivalent $l_{iso}=(l_xl_yl_z)^{1/3}$ is positive $0\%\to29.6\%$ in the 35-case sweep, with no negative values, enabling tailoring of stiffness and frequency via controlled anisotropy and orientation.
